\chapter{Teoreticko-metodologická východiska}
Tato kapitola poskytuje teoretický a metodologický rámec nezbytný pro analýzu a redesign uživatelského rozhraní aplikace. Jejím cílem je představit základní pojmy, principy a trendy v oblasti workplace management software, UX a UI designu, design systémů a evaluace použitelnosti. Součástí kapitoly je rovněž přehled metod a nástrojů, které slouží k identifikaci problémů současného rozhraní a k návrhu optimálního, uživatelsky přívětivého řešení. Teoretická východiska poskytují nezbytný kontext pro praktickou část práce a umožňují aplikovat osvědčené postupy při systematickém redesignu aplikace.

\section{Workplace management software}
% Tuto podkapitolu jsem pro toto zadání přeskočil, ale ve finální diplomové práci jí chci mít přibližně na tomto místě.

\section{UX a UI design}
V oblasti digitálních produktů se pojmy UX (User Experience/Uživatelský zážitek) design a UI (User Interface/Uživatelské rozhraní) design stal klíčovým pro vytváření efektivních a intuitivních rozhraní. UX design je disciplína zaměřující se na řízení celkové zkušenost uživatele při používání nějaké aplikace. UI design je disciplína zaměřující se na vizuální stránku prvků rozhraní.

\subsection{Principy UX designu}
Hlavním cílem UX designu je zajistit, aby uživatelé mohli efektivně, intuitivně a bez zbytečné mentální námahy dosáhnout svých cílů. Podle Dona Normana (Norman, 2013) lze klíčové pricipy UX designu shrnout následovně
\begin{itemize}
    \item Orientace na uživatele (user-centered design) – návrh je založen na potřebách uživatelů
    \item Použitelnost (usability) – ovládání produktu má být snadno naučitelné a jednoduché na používání
    \item Předvídatelnost – prvky a interakce napříč aplikací musí mít jednotný vzhled a chování
    \item Viditelná zpětná vazba – informování uživatele o provedené změně, úspěchu či chybě
    \item Minimalizace kognitivní zátěže – rozhraní by mělo být přehledné, intuitivní a logicky strukturované
\end{itemize}


\subsection{Principy UI designu}
UI design se zaměřuje na vizuální a interakční stránku rozhraní.
\begin{itemize}
    \item Jasná vizuální hierarchie – klíčové prvky jsou zvýrazněny a vedou uživatele logickým způsobem
    \item Konzistence a standardizace – opakující se prvky mají stejný vzhled a stejné chování
    \item Minimalismus a jednoduchost – odstranit nepotřebné prvky, aby rozhraní bylo přehledné
    \item Interakční zpětná vazba – animace, změny barev nebo stavů indikují, že systém reaguje na akce uživatele
\end{itemize}
UI design v praxi úzce souvisí s UX, protože vizuální a interakční prvky ovlivňují, jak snadno a příjemně uživatelé dosahují svých cílů. (Dannaway, 2024)

\subsection{Heuristické zásady}
Heuristická evaluace představuje expertní metodu posuzování rozhraní podle osvědčených principů použitelnosti. Jacob Nielsen (1994) definoval deset základních heuristik, které se často používají při hodnocení a redesignu komplexních aplikací.
\begin{enumerate}
    \item Viditelnost stavu systému
    \item Shoda mezi systémem a reálným světem
    \item Uživatelská kontrola a svoboda
    \item Konzistence a standardy
    \item Prevence chyb
    \item Rozpoznání místo vybavování
    \item Flexibilita a efektivita použití
    \item Estetický a minimalistický design
    \item Pomoc při chybách
    \item Nápověda a dokumentace
\end{enumerate}
Použití heuristických principů je při redesignu stávajících modulů aplikace zásadní pro odhalení nekonzistentních prvků, problémové navigace a oblastí s vysokou kognitivní zátěží.

\section{Design systémy}
Design systémy představují v současnosti jeden z nejdůležitějších nástrojů při návrhu a vývoji digitálních produktů. Reagují na potřebu dlouhodobé konzistence, škálovatelnosti a efektivní spolupráce mezi designéry a vývojáři, zejména v případech rozsáhlých nebo dynamicky se rozvíjejících aplikací. V situaci, kdy produkty obsahují desítky až stovky obrazovek či modulů, se design systém stává klíčovým prvkem umožňujícím udržet jednotnou vizuální a interakční logiku. (Kholmatova, 2017)

\subsection{Definice a účel design systému}
Podle Ally Kholmatové neexistuje žádná jednotná, univerzálně přijímaná definice design systému. V knize Design Systems definuje design systém jako soubor vzájemně propojených vzorů a sdílených postupů, které jsou koherentně uspořádány za účelem podpory digitálních produktů.
Vzory jsou opakující se prvky, které kombinujeme při tvorbě rozhraní. Jde například o uživatelské toky, interakce, tlačítka, formulářová pole, ikony, barvy, typografie nebo microcopy (krátké texty).
Postupy pak znamenají to, jakým způsobem budeme tyto vzory vytvářet, zaznamenávat, sdílet, a používat. Zejména při týmové spolupráci. (Kholmatova, 2017)
Design systém podle zmíněné definice spojuje tři základní roviny
\begin{itemize}
    \item Vizuální prvky – definice estetické podoby zobrazení
    \item Komponentový systém – katalog prvků určených k opakovanému použití
    \item Pravidla a procesy – metodika, jak se jednotlivé prvky používají, kombinují a rozvíjejí
\end{itemize}
Primárním účelem design systému je vytvořit jednotnou a konzistentní uživatelskou zkušenost napříč aplikací, modulární strukturu podporující efektivní vývoj a snížení chybovosti způsobené nesourodostí stylů či komponent.

\subsection{Principy design systémů}
Jedním z nejrozšířenějších přístupů k tvorbě design systémů je Atomic Design, který představil Brad Frost (2016). Tento model dělí uživatelské rozhraní do pěti úrovní
\begin{enumerate}
    \item Atomy – základní prvky jako barvy, typografické styly, ikony či jednoduché UI prvky
    \item Molekuly – kombinace atomů (např. input + ikona)
    \item Organismy – větší komponenty (např. hlavička stránky)
    \item Šablony – strukturální kostra obrazovek
    \item Stránky – konkrétní realizace šablon s reálným obsahem
\end{enumerate}
Tento přístup usnadňuje škálování designu, protože změna na nízké úrovni (např. úprava barevného tokenu) se může automaticky promítnout v celém systému.

\subsection{Přínosy design systémů}
Podle studií společnosti Sparkbox, která prováděla průzkumy Design Systems Survey každý rok mezi lety 2018 až 2022, dokáže organizace s funkčním design systémem zkrátit čas tvorby obrazovek až o 50 %. Další výhody zahrnují
\begin{itemize}
    \item Rychlejší vývoj – vývojáři opakovaně používají již schválené komponenty místo vytváření nových
    \item Snížení technického dluhu – díky jednotné knihovně komponent se eliminuje nekonzistence a potřeba předělávek
    \item Lepší komunikace v týmu – designéři, vývojáři a produktoví manažeři pracují se stejným pojmovým rámcem
    \item Snazší onboarding nových členů – noví designéři a vývojáři se rychleji orientují ve vizuálním jazyce produktu
    \item Vyšší kvalita UX – repetitivní používání testovaných komponent vede k menší chybovosti uživatelů a vyšší spokojenosti (Sparkbox, n. r.)
\end{itemize}
    

% \subsection{Design systém}
% Software musí splňovat vyšší nároky na profesionální vizuální identitu, intuitivnost a očekávání zákazníků. Dhanwani (2025) uvádí, že moderní B2B produkty jsou hodnoceny nejen podle funkcionality, ale především podle kvality uživatelské zkušenosti – a design systém je jedním z hlavních nástrojů, jak tuto kvalitu zajistit.
% Design systém umožní
% •	efektivní přidávání nových modulů bez ztráty konzistence
% •	snazší přizpůsobení produktu specifickým požadavkům klientů
% •	budování vizuální identity podporující důvěru v produkt
% •	dlouhodobě udržitelný rozvoj bez návratu k chaotickému organickému růstu

\section{Metody použitelnosti a evaluace}
Uživatelský výzkum, testování použitelnosti, vhodné metody